% Nouveau EGA V (cote 25): archive cote 25, site page 161 (PDF page 162).
% typescript. Draft compared: none.
% AI-transcribed by Claude Opus 5.5 (claude-opus-5-5, adaptive thinking, effort high), Message Batches API, batch ega25b (msgbatch_01A7n5YuXubUPNK1zsepi8RG), 2026-11-13; not yet checked by eye.
\scanpage{25/161}
\mspage{4}

puissance de chaque v\xout{X}ariable $X_i$ . Plus précis\unsure{\xout{x}}ément, le sous-\xout{pré}schéma
de $P^r$ défini par l'idéal homogène engendré par \xout{le} $\phi$ et les $\frac{\partial\phi}{\partial X_i}$ est
exactement formé des points de V($\phi$) en lesquels V($\phi$) n'est pas lisse sur
Spec(A) avec une dimension relative r-1.
15.2. Soi\unsure{e}nt X un préschéma, $\underline{J}$ un Idéal \inserted{qu coh} \xout{dexxxx} sur X, nous dirons que
l'homomorphisme ``d'augmentation'' $\underline{O}_X \to \underline{O}_X/\underline{J}$ est une ``augmentation
élémentaire de multiplicité n'' si elle satisfait aux conditions suivantes
a) $\underline{J}/\underline{J}^2 = \underline{N}_{Y/X}$ est loc libre de type fini sur Y = V($\underline{J}$).
b) Le noyau de l'homomorphisme canonique
\[ \hand{\mu} : \underline{Sym}_{\underline{O}_Y}(\underline{N}) \longrightarrow gr_{\underline{J}}(\underline{O}_X) \]
est engendré par le noyau \inserted{K} de $\hand{\mu}^n$, lequel est un sous-Module inversible
localement facteur direct de $\underline{Sym}^n_{\underline{O}_Y}(\underline{N})$ (i.e. est engendrable localement
sur Y par une section $\phi$ qui ne\xout{xxx} s'annule en aucun point).
c) Ladite $\phi$ (qui n'est définie, localement, que modulo multiplication
par une unité) est une forme \uline{lisse} . (NB On aurait pu introduire dans
\margin{\textbar}15.1. la notion de sous-Module inversible lisse de $\underline{Sym}^n(E)$, à vrai dire
géométriquement plus importante que celle de \uline{section} lisse, car V($\phi$) est
en fait défini par un tel sous-Module ...\xout{;} cela aurait l'avantage de
permettre de formuler b) et c) en une seule condition).

Lorsque n=2, on parlera de ``augmentation \inserted{de type} quadratique élémentaire'' ,
Si A est un anneau, \xout{lxxxlxxxxxthéxiexx} muni d'une \unsure{augmentation|augmantation} A$\to$A/J,
on conviendra encore de dire que cette augmentation est élémentaire de
multiplicité \unsure{n} , s'il en est ainsi pour .... \xout{E}Lorsque A est un anneau
local, on dira simplement par abus de langage que A est ``élémentaire de
multiplic\uline{i}té n'' si l'aug\xout{m}mentation A $\to$ A/$\underline{r}$(A) est élémentaire de
multiplicité \unsure{n}; noter que cela implique $n \geqslant 2$ , et si A est noethérien,
A est nécessairement non régulier. Si X est un préschéma, n un entier

\mspage{161}
